En una actividad de dibujo técnico en el colegio, un estudiante ubica tres puntos en el plano que corresponden a las posiciones de tres postes en el patio. Al intentar unirlos con una sola línea recta, observa que los puntos no quedan alineados.

Sin embargo, el estudiante afirma que es posible trazar una única recta que pase exactamente por los tres puntos, ajustando correctamente la inclinación de la regla.

¿La afirmación del estudiante es verdadera o falsa?

Justifique su respuesta utilizando argumentos geométricos claros y explicando con sus propias palabras la condición necesaria para que varios puntos pertenezcan a una misma recta. No se aceptan respuestas sin justificación.